\section{Introdução}
\label{sec:intro}

Sistemas de reação-difusão produzem, a partir de um estado quase homogêneo, estruturas espaciais
organizadas (pontos, listras, labirintos, espirais e ``planadores'') cuja forma depende de poucos
parâmetros \cite{turing1952,pearson1993}. A pergunta que orienta este trabalho é a mesma para
todos os sistemas físicos estudados pela turma: \emph{onde estão as estruturas de interesse e como
elas evoluem ao longo do tempo?} Responder a ela exige uma segmentação por pixel, e não apenas a
classificação do regime de cada simulação.

O sistema estudado pela equipe é o \texttt{gray\_scott\_reaction\_diffusion} do The Well
\cite{ohana2024well}: 1200 trajetórias do modelo de Gray-Scott \cite{grayscott1984,pearson1993}
em seis regimes de parâmetros, com 1001 instantes de $128\times128$ pixels cada. O conjunto não
traz máscaras de segmentação, e o ground truth (GT) é derivado da concentração $B$ por um critério
físico documentado e congelado antes da comparação (Seção~\ref{sec:gt}).

O experimento segue o protocolo comum da turma, que permite juntar os resultados de 11 sistemas
físicos 2D num único benchmark:
\begin{itemize}
  \item resolução com maior eixo de no máximo 128 px e sem ampliação;
  \item no máximo 25 GB de dados por sistema;
  \item três arquiteturas: U-Net \cite{ronneberger2015}, SegFormer \cite{xie2021segformer} e
    Mask2Former \cite{cheng2022mask2former};
  \item 20 épocas, com o checkpoint da melhor época de validação;
  \item dez pré-processamentos (P1--P10) e cinco pós-processamentos (Q1--Q5), cada um testado
    isoladamente contra o baseline.
\end{itemize}

São avaliadas quatro condições:
\begin{enumerate}
  \item sem pré e sem pós (referência);
  \item somente pré (efeito isolado de P1--P10);
  \item somente pós (efeito isolado de Q1--Q5);
  \item pré e pós (as $10\times3\times5 = 150$ combinações).
\end{enumerate}

As contribuições são:
\begin{itemize}
  \item um protocolo reprodutível de segmentação do Gray-Scott, com GT físico, entrada restrita a
    $A$ para evitar vazamento de rótulo, divisão oficial por trajetória e operações que respeitam o
    contorno periódico;
  \item a análise de \emph{por que} cada pré e pós-processamento melhora ou piora a segmentação,
    relacionando o efeito ao mecanismo de cada arquitetura (convolução com \emph{skip connections},
    atenção hierárquica em 1/4 da resolução e classificação de máscaras por consultas) e à física
    de cada regime;
  \item tabelas de resultados no esquema padronizado do protocolo (CSV e \LaTeX), prontas para a
    junção com os demais sistemas no artigo de benchmark.
\end{itemize}
